% ============================================================
% exemples.tex
% Fichier de démonstration du package outilsprof
% Auteur : Jean Roussie
% Date : 2026/09/11
% ============================================================

\documentclass[11pt,a4paper]{article}

% --- Encodage et langue ---
\usepackage[utf8]{inputenc}
\usepackage[T1]{fontenc}
\usepackage[french]{babel}
\usepackage{lmodern}

% --- Mise en page ---
\usepackage[margin=2cm]{geometry}
\usepackage{parskip}

% --- Mathématiques ---
\usepackage{amsmath,amssymb}

% --- Le package à démontrer ---
\usepackage{outilsprof}

% --- Titre ---
\title{\textbf{Démonstration du package \texttt{outilsprof}}\\
       \large Outils pour l'enseignement des mathématiques}
\author{Jean Roussie \\ \small Collège La Boétie — Sarlat (Dordogne)}
\date{\today}

\begin{document}

\maketitle

\tableofcontents
\newpage

% ============================================================
\section{Axe gradué}
% ============================================================

\subsection{Axe par défaut}

\begin{center}
\axegradue
\end{center}

\subsection{Axe de 0 à 400, pas de 100, 5 sous-intervalles}

\begin{center}
\axegradue[longueur=0.8, fin=400, pas=100, nombre=5]
\end{center}

\subsection{Axe de 100 à 500}

\begin{center}
\axegradue[debut=100, fin=500, pas=100, nombre=5]
\end{center}

\subsection{Axe avec points}

\begin{center}
\axegradue[fin=400, pas=100, nombre=5, points={A/120, B/340/blue, C/60/green}]
\end{center}

\subsection{Variante \texttt{\textbackslash axegraduepoints}}

\begin{center}
\axegraduepoints[longueur=0.8, fin=400, pas=100, nombre=5]{A/120, B/340/blue, C/60/green}
\end{center}

% ============================================================
\section{Repère gradué}
% ============================================================

\subsection{Repère par défaut}

\begin{center}
\reperegradue
\end{center}

\subsection{Fenêtre personnalisée $[-5,5]\times[-5,5]$, grille de 1}

\begin{center}
\reperegradue[bg={(-5,-5)}, hd={(5,5)}, grille=1]
\end{center}

\subsection{Fenêtre $[-4,3]\times[-2,6]$, grille de 0.5}

\begin{center}
\reperegradue[bg={(-4,-2)}, hd={(3,6)}, grille=0.5]
\end{center}

\subsection{Grille désactivée}

\begin{center}
\reperegradue[bg={(-5,-5)}, hd={(5,5)}, grille=0]
\end{center}

\subsection{Graduations principales personnalisées}

\subsubsection{pasx=2, pasy=2}

\begin{center}
\reperegradue[bg={(-10,-10)}, hd={(10,10)}, pasx=2, pasy=2, grille=1]
\end{center}

\subsubsection{pasx=5, pasy=5, nombrex=5, nombrey=5}

\begin{center}
\reperegradue[bg={(-10,-10)}, hd={(10,10)},
              pasx=5, pasy=5, nombrex=5, nombrey=5, grille=1]
\end{center}

% ============================================================
\section{Étiquette de l'intersection des axes}
% ============================================================

\subsection{Origine en $O(0,0)$}

\begin{center}
\reperegradue[bg={(-5,-5)}, hd={(5,5)}]
\end{center}

\subsection{Intersection en $\Omega(2;3)$}

\begin{center}
\reperegradue[bg={(-5,-5)}, hd={(5,5)}, origine={(2,3)}]
\end{center}

\subsection{Intersection en $\Omega(-3;-2)$}

\begin{center}
\reperegradue[bg={(-5,-5)}, hd={(5,5)}, origine={(-3,-2)}]
\end{center}

\subsection{Intersection sur l'axe des ordonnées : $\Omega(0;3)$}

\begin{center}
\reperegradue[bg={(-5,-5)}, hd={(5,5)}, origine={(0,3)}]
\end{center}

% ============================================================
\section{Placement de points}
% ============================================================

\subsection{Un seul point (rouge par défaut)}

\begin{center}
\reperegradue[bg={(-6,-6)}, hd={(6,6)}, points={A/(4,5)}]
\end{center}

\subsection{Points avec couleurs}

\begin{center}
\reperegradue[bg={(-6,-6)}, hd={(6,6)}, grille=1,
              points={A/(4,5), B/(-2,3)/blue, C/(1,-2)/green}]
\end{center}

\subsection{Couleurs composées TikZ}

\begin{center}
\reperegradue[bg={(-4,-4)}, hd={(4,4)}, grille=0.5,
              points={D/(0.5,-3.75)/orange!80!black,
                      E/(-1.5,2.25)/red!70}]
\end{center}

\subsection{Points sur les axes et à l'origine}

\begin{center}
\reperegradue[bg={(-5,-5)}, hd={(5,5)},
              points={O/(0,0)/black, X/(5,0)/blue, Y/(0,5)/green}]
\end{center}

% ============================================================
\section{Options d'affichage des points}
% ============================================================

\subsection{Noms affichés, coordonnées masquées (défaut)}

\begin{center}
\reperegradue[bg={(-5,-5)}, hd={(5,5)},
              points={A/(4,5), B/(-2,3)/blue},
              affichernom=true, affichercoords=false]
\end{center}

\subsection{Noms et coordonnées affichés}

\begin{center}
\reperegradue[bg={(-5,-5)}, hd={(5,5)},
              points={A/(4,5), B/(-2,3)/blue},
              affichernom=true, affichercoords=true]
\end{center}

\subsection{Noms masqués, coordonnées affichées}

\begin{center}
\reperegradue[bg={(-5,-5)}, hd={(5,5)},
              points={A/(4,5), B/(-2,3)/blue},
              affichernom=false, affichercoords=true]
\end{center}

\subsection{Noms et coordonnées masqués (points seuls)}

\begin{center}
\reperegradue[bg={(-5,-5)}, hd={(5,5)},
              points={A/(4,5), B/(-2,3)/blue},
              affichernom=false, affichercoords=false]
\end{center}

% ============================================================
\section{Variante \texttt{\textbackslash reperegraduepoints}}
% ============================================================

\begin{center}
\reperegraduepoints[bg={(-5,-5)}, hd={(5,5)}]{A/(2,3), B/(-1,-2)/orange}
\end{center}

% ============================================================
\section{Cas combiné : toutes les options}
% ============================================================

\begin{verbatim}
\reperegradue[
    bg=(-4,-3), hd=(6,4),
    origine=(0,0),
    grille=0.5,
    pasx=1, pasy=1,
    nombrex=2, nombrey=2,
    points={A/(3,2), B/(-2,1)/blue, C/(1,-2)/green},
    affichernom=true, affichercoords=true
]
\end{verbatim}

\begin{center}
\reperegradue[
    bg={(-4,-3)}, hd={(6,4)},
    origine={(0,0)},
    grille=0.5,
    pasx=1, pasy=1,
    nombrex=2, nombrey=2,
    points={A/(3,2), B/(-2,1)/blue, C/(1,-2)/green},
    affichernom=true, affichercoords=true
]
\end{center}

% ============================================================
\section{Utilisation pédagogique}
% ============================================================

\subsection{Placer des points dans un repère}

\begin{methode}
    Pour placer un point $A(x_A, y_A)$ dans un repère :
    \begin{enumerate}
        \item Repérer $x_A$ sur l'axe des abscisses ;
        \item Repérer $y_A$ sur l'axe des ordonnées ;
        \item Tracer les parallèles aux axes passant par ces repères ;
        \item Le point $A$ est à l'intersection.
    \end{enumerate}
\end{methode}

\begin{center}
\reperegradue[bg={(-3,-3)}, hd={(5,5)},
              points={A/(3,2)/blue, B/(-2,1)/green, C/(1,-2)/red}]
\end{center}

\subsection{Lire les coordonnées d'un point}

\begin{retenir}
    Un point $M$ du plan est repéré par un couple unique $(x_M, y_M)$ :
    $x_M$ est son abscisse, $y_M$ son ordonnée.
\end{retenir}

\begin{center}
\reperegradue[bg={(-4,-4)}, hd={(4,4)},
              points={M/(2,3)/orange}]
\end{center}

\begin{vigilance}
    Ne pas confondre \textbf{abscisse} (axe horizontal) et
    \textbf{ordonnée} (axe vertical) !
\end{vigilance}

% ============================================================
\section{Boîtes colorées}
% ============================================================

\subsection{Boîte « À retenir » (bleu par défaut)}

\begin{retenir}
    Le théorème de Pythagore : dans un triangle rectangle, le carré de
    l'hypoténuse est égal à la somme des carrés des deux autres côtés.
    \[
        a^2 + b^2 = c^2
    \]
\end{retenir}

\subsection{Boîte « Attention ! » (rouge par défaut)}

\begin{vigilance}
    Ne pas oublier de vérifier les unités avant tout calcul !
    Un résultat sans unité est un résultat faux.
\end{vigilance}

\subsection{Boîte « Méthode » (vert par défaut)}

\begin{methode}
    Pour résoudre une équation du premier degré :
    \begin{enumerate}
        \item Regrouper les termes en $x$ dans un membre ;
        \item Regrouper les constantes dans l'autre membre ;
        \item Simplifier chaque membre ;
        \item Diviser par le coefficient de $x$.
    \end{enumerate}
\end{methode}

\subsection{Boîtes avec couleur personnalisée}

\begin{retenir}[orange]
    Version personnalisée en orange : utile pour mettre en avant une
    notion secondaire.
\end{retenir}

\begin{vigilance}[purple]
    Version personnalisée en violet : pour une remarque intermédiaire.
\end{vigilance}

\begin{methode}[cyan!80!black]
    Version personnalisée en cyan foncé.
\end{methode}

\subsection{Versions sous forme de commandes}

\boiteRetenir{Version courte avec la commande \texttt{\textbackslash boiteRetenir}.}

\boiteVigilance{Version courte avec la commande \texttt{\textbackslash boiteVigilance}.}

\boiteMethode{Version courte avec la commande \texttt{\textbackslash boiteMethode}.}

% ============================================================
\section{Boîte de calcul}
% ============================================================

\subsection{Boîte de calcul par défaut (10\,cm $\times$ 3\,cm)}

\begin{center}
\boitecalcul{$2 + 2 = 4$}
\end{center}

\subsection{Boîte de calcul personnalisée (8\,cm $\times$ 4\,cm)}

\begin{center}
\boitecalcul[8cm][4cm]{$\dfrac{3}{4} + \dfrac{1}{2} = \dfrac{5}{4}$}
\end{center}

\subsection{Boîte de calcul large (12\,cm $\times$ 5\,cm)}

\begin{center}
\boitecalcul[12cm][5cm]{%
    $\displaystyle \int_0^1 x^2 \, \mathrm{d}x = \frac{1}{3}$%
}
\end{center}

% ============================================================
\section{Combinaison des outils}
% ============================================================

\begin{methode}
    Pour placer un point $A$ d'abscisse 120 sur un axe gradué, on
    repère la position entre 100 et 200, puis on utilise les graduations
    secondaires.
\end{methode}

\begin{center}
\axegradue[fin=400, pas=100, nombre=5, points={A/120/red, B/340/blue}]
\end{center}

\begin{retenir}
    Un point est entièrement déterminé par son abscisse sur l'axe.
\end{retenir}

\begin{vigilance}
    Attention au sens de lecture : l'axe est orienté de la gauche vers
    la droite.
\end{vigilance}

\end{document}